% Part: sets-functions-relations
% Chapter: infinite
% Section: dedekind
%
\documentclass[../../../include/open-logic-section]{subfiles}

\begin{document}

\olfileid[ar]{sfr}{infinite}{dedekind}
\olsection{جبرات ديديكند}

لا نريد للأعداد الطبيعية أن تكون لا نهائية فحسب؛ بل نريد لها أن تتمتع
بخواص (جبرية) معينة: فيلزم أن تسلك سلوكًا حسنًا تحت الجمع والضرب وما
إلى ذلك.

كانت فكرة ديديكند أن يتخذ مفهوم \emph{دالة اللاحق} مفهومًا أساسيًا، ثم
يصف الأعداد بأنها الأشياء ذات الخواص الآتية:
\begin{enumerate}
	\item يوجد عدد، $0$، ليس لاحقًا لأي عدد
	\\أي إن $0 \notin \ran{s}$
	\\أي إن $\forall x\ s(x) \neq 0$
	\item للأعداد المتمايزة لواحق متمايزة
	\\أي إن $s$ هي !!a{injection}
	\\أي إن $\forall x \forall y (s(x) = s(y) \lif x = y)$
	\item\ollabel{repeatedapplication} يُحصل على كل عدد من
	$0$ بتطبيق دالة اللاحق مرارًا.
\end{enumerate}
يسهل تناول الشرطين الأولين باستعمال منطق الرتبة الأولى (انظر أعلاه).
لكن لا يمكننا تناول \olref{repeatedapplication} باستعمال منطق الرتبة
الأولى وحده. وكان إنجاز ديديكند الحاسم أن أعاد صياغة الشرط
\olref{repeatedapplication} بلغة نظرية المجموعات كما يأتي:
\begin{enumerate}
	\item[3$'$.] الأعداد الطبيعية هي أصغر مجموعة
	\emph{منغلقة تحت دالة اللاحق}: أي إننا إذا طبقنا $s$ على أي
	!!{element} من المجموعة، حصلنا على !!{element} آخر من المجموعة.
\end{enumerate}
لكن سيلزمنا أن نبسط ذلك على مهل.

\begin{defn}\ollabel{Closure}
	لكل دالة $f$، تكون المجموعة $X$ \emph{منغلقة} تحت $f$
	{إذا وفقط إذا} كان $(\forall x \in X)f(x) \in X$. والآن نعرّف، لكل $o$:
	$$\closureofunder{f}{o} = \bigcap\Setabs{X}{o \in X\text{ و}X
\text{ منغلقة تحت $f$}}$$
\end{defn}

إذن $\closureofunder{f}{o}$ هي تقاطع جميع المجموعات المنغلقة تحت $f$
التي $o$ فيها !!a{element}. ومن ثم، حدسيًا، تكون
$\closureofunder{f}{o}$ \emph{أصغر} مجموعة منغلقة تحت $f$، و$o$ فيها
!!a{element}. وتجعل النتيجة الآتية هذه الفكرة الحدسية دقيقة؛
\begin{lem}\ollabel{closureproperties}
	لكل دالة $f$ ولكل $o \in A$:
	\begin{enumerate}
		\item\ollabel{closurehaselem} $o \in \closureofunder{f}{o}$؛ و
		\item\ollabel{closureclosed} $\closureofunder{f}{o}$ منغلقة تحت $f$؛ و
		\item\ollabel{closuresmallest} إذا كانت $X$ منغلقة تحت $f$ وكان $o \in
		X$، فإن $\closureofunder{f}{o} \subseteq X$
	\end{enumerate}
\end{lem}

\begin{proof}
لاحظ أن ثمة مجموعة واحدة على الأقل منغلقة تحت $f$، و$o$ فيها
!!a{element}، وهي $\ran{f}\cup
\{o\}$. ولذلك يوجد $\closureofunder{f}{o}$، أي تقاطع
\emph{جميع} المجموعات من هذا القبيل. ويلزمنا الآن أن نتحقق من
\olref{closurehaselem}--\olref{closuresmallest}.

فيما يخص \olref{closurehaselem}: لدينا $o \in
\closureofunder{f}{o}$، لأنها تقاطع مجموعات، و$o$ في كل منها
!!a{element}.

فيما يخص \olref{closureclosed}: افترض أن $x \in
\closureofunder{f}{o}$. فإذا كان $o \in X$ وكانت $X$ منغلقة تحت $f$،
كان $x \in X$، وعندئذ يكون $f(x) \in X$ لأن $X$ منغلقة تحت $f$.
ومن ثم $f(x) \in \closureofunder{f}{o}$.

فيما يخص \olref{closuresmallest}: بوجه عام، إذا كان $X
\in C$ فإن $\bigcap C \subseteq X$.
\end{proof}

وباستخدام ذلك يمكننا أن نقول:

\begin{defn}
\emph{جبر ديديكند} هو مجموعة $A$ مع دالة $f
\colon A \to A$ وعنصر ما $o \in A$ بحيث:
	\begin{enumerate}
		\item \ollabel{ded:proper} $o \notin \ran{f}$
		\item \ollabel{ded:injection} $f$ هي !!a{injection}
		\item \ollabel{ded:closure} $A = \closureofunder{f}{o}$
	\end{enumerate}
\end{defn}

بما أن $A = \closureofunder{f}{o}$، تخبرنا نتيجتنا السابقة أن $A$ أصغر
مجموعة منغلقة تحت $f$، و$o$ فيها !!a{element}. ومن الواضح أن جبر
ديديكند لانهائي بحسب ديديكند؛ فما عليك إلا أن تنظر إلى الشرطين
\olref{ded:proper} و\olref{ded:injection} من التعريف. لكن الحقيقة الأكثر
إثارة هي أن كل مجموعة لانهائية بحسب ديديكند يمكن تحويلها إلى جبر
ديديكند.

\begin{thm}\ollabel{thm:DedekindInfiniteAlgebra}
إذا وجدت مجموعة لانهائية بحسب ديديكند، وجد جبر ديديكند.
\end{thm}

\begin{proof}
لتكن $D$ لانهائية بحسب ديديكند. إذن توجد دالة متباينة $g \colon D \to
D$ وعنصر $o  \in D \setminus \ran{g}$. والآن لنجعل $A =
\closureofunder{g}{o}$؛ وبحسب \olref{closureproperties}، توجد $A$ ويكون
$o \in A$. ولنجعل $f =
\funrestrictionto{g}{A}$. وسنبين أن $A, f, o$ تشكل جبر ديديكند.

فيما يخص \olref{ded:proper}: لدينا $o \notin \ran{g}$ و$\ran{f}
\subseteq \ran{g}$، ومن ثم $o\notin \ran{f}$.

فيما يخص \olref{ded:injection}: الدالة $g$ متباينة على $D$؛ ولذلك يجب
أن تكون $f \subseteq g$ متباينة.

فيما يخص \olref{ded:closure}: بحسب \olref{closureproperties}، تكون $A$
منغلقة تحت $g$؛ ومن باب أولى تكون $A$ منغلقة تحت $f$. ومن ثم
$\closureofunder{f}{o} \subseteq A$ بحسب \olref{closureproperties}.
وبما أن $\closureofunder{f}{o}$ منغلقة أيضًا تحت $f$ وأن $f =
\funrestrictionto{g}{A}$، يلزم أن تكون $\closureofunder{f}{o}$ منغلقة
تحت $g$. ولذلك $A \subseteq \closureofunder{f}{o}$ بحسب
\olref{closureproperties}.
\end{proof}

\end{document}
